% vim: set textwidth=78 autoindent:

\subsection{Otros complementos del núcleo}

% when the revision of a section has been finalized, 
% comment out the following line:
% \updatedisclaimer

Los otros complementos del núcleo son listados en la tabla \ref{tab:other_core}, junto con referencias a los  
capítulos en este manual los cuales cubren su uso.

\begin{table}[H]
\centering
\caption{Otros complementos del núcleo}\label{tab:other_core}\medskip
\small
 \begin{tabular}{|l|l|p{4in}|}
\hline \textbf{Ícono} & \textbf{Complemento} & \textbf{Manual de Referencia}\\
\hline
\includegraphics[width=0.6cm]{diagram_overlay}
 & Sobreposición de diagramas \index{plugins!diagram}& Capítulo \ref{sec:diagram}\\
\hline
\includegraphics[width=0.6cm]{grass}
 & GRASS \index{plugin!grass toolbox} & Capítulo \ref{sec:grass} and Appendix \ref{appdx_grass_toolbox_modules}\\
 \hline
\includegraphics[width=0.6cm]{plugin_installer}
 & Instalador de Complementos \index{plugins!Plugin Installer} & Capítulo \ref{sec:python_plugin_installer}\\
\hline
\includegraphics[width=0.6cm]{spiticon}
 & SPIT \index{plugins!spit}& Capítulo \ref{sec:loading_postgis_data} \\
 \hline
\includegraphics[width=0.6cm]{mIconAddWfsLayer}
 & WFS & Capítulo \ref{sec:ogc-wfs} \\
\hline
\end{tabular}
\end{table}
